\section{Design a Logging Framework}
\label{sec:case-logger}

\problemstatement{Design a logging framework supporting severity levels
(DEBUG, INFO, WARN, ERROR), multiple output sinks (console, file,
network), a configurable minimum level, and the ability to add sinks or
formats without changing call sites.}

\begin{patternbox}[Clarify Before Designing]
Ask: What levels and can they be filtered globally or per sink? Multiple
simultaneous sinks? Thread-safe? Should a message be passed along a chain
of level handlers or broadcast to all sinks? For this design assume:
standard levels, multiple sinks, a global threshold, and thread safety as
a discussion point.
\end{patternbox}

\subsection*{Identify the entities}

A \texttt{Logger} (front door, often a \textbf{Singleton}), a
\texttt{LogLevel} enum, \texttt{LogSink}/\texttt{Appender} implementations
(\texttt{ConsoleSink}, \texttt{FileSink}), and a \texttt{Formatter}. Two
classic pattern framings apply: broadcast to all sinks (\textbf{Observer})
or pass up a \textbf{Chain of Responsibility} of level handlers.

\begin{ideabox}[Level Filtering + Pluggable Sinks]
Keep a single global threshold so sub-threshold messages are dropped
cheaply. Model each output as a \texttt{LogSink} plug-in the logger writes
to -- adding a sink never touches call sites. A \textbf{Chain of
Responsibility} variant instead routes a message through DEBUG$\to$INFO
$\to$WARN$\to$ERROR handlers, each deciding whether to act; both are valid,
name the trade-off.
\end{ideabox}

\subsection*{Class design (C++ skeleton)}

\begin{lstlisting}
#include <bits/stdc++.h>
using namespace std;

enum class Level { Debug, Info, Warn, Error };

struct LogSink {                                // pluggable output
    virtual void write(Level lvl, const string& msg) = 0;
    virtual ~LogSink() = default;
};
struct ConsoleSink : LogSink {
    void write(Level, const string& msg) override { cout << "[console] " << msg << "\n"; }
};
struct FileSink : LogSink {
    void write(Level, const string& msg) override { cout << "[file] " << msg << "\n"; }
};

class Logger {                                  // Singleton front door
    Level threshold = Level::Info;
    vector<unique_ptr<LogSink>> sinks;
    Logger() = default;
public:
    static Logger& instance() { static Logger l; return l; }
    void setThreshold(Level l) { threshold = l; }
    void addSink(unique_ptr<LogSink> s) { sinks.push_back(move(s)); }

    void log(Level lvl, const string& msg) {
        if (lvl < threshold) return;            // cheap filter
        for (auto& s : sinks) s->write(lvl, msg);   // broadcast to sinks
    }
};

int main() {
    Logger& log = Logger::instance();
    log.addSink(make_unique<ConsoleSink>());
    log.addSink(make_unique<FileSink>());
    log.setThreshold(Level::Info);
    log.log(Level::Debug, "verbose");           // filtered out
    log.log(Level::Error, "something broke");   // goes to both sinks
    return 0;
}
\end{lstlisting}

\begin{takeawaybox}
A logger design lets you name several patterns naturally: \textbf{Singleton}
for the global logger, pluggable \textbf{Sinks} (Strategy/Observer) for
outputs, and optionally a \textbf{Chain of Responsibility} over levels.
The must-mention non-functional concern is \emph{thread safety} -- a lock
or a lock-free queue around the sink writes at scale.
\end{takeawaybox}
